\section{What the result does and does not show}
The target is fixed and shallow. The adverse dependence comes from the depth of the learner, which allows a low-cost implementation of the observed noisy labels. The data are IID and the noise probability is not selected from the sample. All likelihood temperatures fail on the same event, so knowledge of the true noise variance or data-dependent temperature tuning does not avoid the example.

For tanh, the finite-prior statement covers every fixed prefactor range $0<Ns\le\nu_{\max}<1$, including mean-field scaling with any fixed prefactor below one and every exponent $\gamma>1$ at sufficiently large widths. The constants may deteriorate as $\nu_{\max}$ approaches one. The canonical prefactor $Ns=1$ is outside this theorem: its matched depth cost $\kappa(1)$ is zero, so a different argument is needed there. The ReLU extension has the separately stated $s\le N^{-1-\kappa}$ scope. These are results for the specified Gaussian architectures and observation models.

The sample size and the required learner width are exponential in learner depth, and input dimension grows linearly with that depth. The result rules out a guarantee polynomial jointly in input dimension, learner depth, teacher width, inverse excess accuracy, and confidence uniformly over learner width for this unrestricted Gaussian family. It does not by itself rule out a fixed-dimension, fixed-accuracy theorem as depth grows, or a bound restricted to learner widths polynomial in depth.

These are conventional mathematical proofs. Independent mathematical review records and the exact reviewed sources accompany the manuscript; no Lean verification is claimed.
